\documentclass[12pt,a4paper]{article}

% --- Packages ---
\usepackage[utf8]{inputenc}
\usepackage[T1]{fontenc}
\usepackage[french]{babel}
\usepackage{geometry}
\usepackage{enumitem}
\usepackage{array}
\usepackage{fancyhdr}
\usepackage{amsmath, amssymb}

% --- Mise en page ---
\geometry{margin=2.5cm}

% --- Pied de page ---
\pagestyle{fancy}
\fancyhf{}
\fancyfoot[R]{\it{G.M.}}
\renewcommand{\headrulewidth}{0pt}
\renewcommand{\footrulewidth}{0pt}

\setlength{\parindent}{0pt}  % pas d'indentation pour tous les paragraphes

\begin{document}

\begin{center}
    \Large \textbf{Automatismes - QCM} \\[5pt]
    \normalsize Aucune justification n’est demandée et une seule réponse est possible
par question.
\end{center}

\vspace{0.7cm}

\textbf{Question 1 :} 
\vspace{5pt}

Jean consacre 25\% de sa journée de dimanche à faire ses devoirs.
80\% du temps consacré aux devoirs est consacré à faire un exposé.
\\Le pourcentage du temps consacré à l’exposé par rapport à la journée de dimanche
est égal à :
\vspace{0.3cm}

\renewcommand{\arraystretch}{3} % multiplie la hauteur par 2
\begin{tabular}{|m{3.4cm}|m{3.4cm}|m{3.4cm}|m{3.4cm}|}
\hline
\centering \vspace{-0.3cm} \textbf{A.} $80\% - 25\%$ & 
\centering \vspace{-0.3cm} \textbf{B.} $\dfrac{1}{4} \times 80\%$ & 
\centering \vspace{-0.3cm} \textbf{C.} $0,08 \times 25\%$ & 
\centering \vspace{-0.3cm} \textbf{D.} Cela dépend de la durée de la journée 
du dimanche \tabularnewline
\hline
\end{tabular}

\vspace{0.7cm}

\textbf{Question 1 :} 
\vspace{5pt}

Le prix d'un article est passé de 250 € à 200 €.  
Par conséquent, son prix a été multiplié par : 
\vspace{0.3cm}

\renewcommand{\arraystretch}{3} 
\begin{tabular}{|m{3.4cm}|m{3.4cm}|m{3.4cm}|m{3.4cm}|}
\hline
\centering \vspace{-0.3cm} \textbf{A.} 1,25 & 
\centering \vspace{-0.3cm} \textbf{B.} 0,75 & 
\centering \vspace{-0.3cm} \textbf{C.} 0,8 & 
\centering \vspace{-0.3cm} \textbf{D.} $-0,8$ \tabularnewline
\hline
\end{tabular}

\vspace{0.7cm}

\textbf{Question 2 :} 
\vspace{5pt}

Le prix d'un produit est passé de 80 € à 100 €.  
Il a donc été multiplié par : 
\vspace{0.3cm}

\renewcommand{\arraystretch}{3} 
\begin{tabular}{|m{3.4cm}|m{3.4cm}|m{3.4cm}|m{3.4cm}|}
\hline
\centering \vspace{-0.3cm} \textbf{A.} 0,75 & 
\centering \vspace{-0.3cm} \textbf{B.} 1,25 & 
\centering \vspace{-0.3cm} \textbf{C.} 1,75 & 
\centering \vspace{-0.3cm} \textbf{D.} $\dfrac{25}{100}$ \tabularnewline
\hline
\end{tabular}

%-------------------------
\vspace{0.7cm}

\textbf{Question 3 :} 
\vspace{5pt}

Dans une classe, 12 élèves pratiquent le basket et représentent 30\% des élèves de la classe.  
Le nombre total d'élèves de cette classe est : 
\vspace{0.3cm}

\renewcommand{\arraystretch}{3} 
\begin{tabular}{|m{3.4cm}|m{3.4cm}|m{3.4cm}|m{3.4cm}|}
\hline
\centering \vspace{-0.3cm} \textbf{A.} 33 & 
\centering \vspace{-0.3cm} \textbf{B.} 30 & 
\centering \vspace{-0.3cm} \textbf{C.} 40 & 
\centering \vspace{-0.3cm} \textbf{D.} 30/12 \tabularnewline
\hline
\end{tabular}

\vspace{0.7cm}

\textbf{Question 4 :} 
\vspace{5pt}

Sur 200 habitants, 50 sont des enfants.  
La proportion d'enfants dans la population est : 
\vspace{0.3cm}

\renewcommand{\arraystretch}{3} 
\begin{tabular}{|m{3.4cm}|m{3.4cm}|m{3.4cm}|m{3.4cm}|}
\hline
\centering \vspace{-0.3cm} \textbf{A.} 0,25 & 
\centering \vspace{-0.3cm} \textbf{B.} 50 & 
\centering \vspace{-0.3cm} \textbf{C.} 25 & 
\centering \vspace{-0.3cm} \textbf{D.} 1,25 \tabularnewline
\hline
\end{tabular}

%-------------------------
\vspace{0.7cm}
\newpage

\textbf{Question 5 :} 
\vspace{5pt}

Un produit vaut 200 € et son prix diminue de 10\%.  
Son nouveau prix est : 
\vspace{0.3cm}

\renewcommand{\arraystretch}{3} 
\begin{tabular}{|m{3cm}|m{3.4cm}|m{3.8cm}|m{3.4cm}|}
\hline
\centering \vspace{-0.3cm} \textbf{A.} $200 \times 0,1$ & 
\centering \vspace{-0.3cm} \textbf{B.} $200 \times \left(-\dfrac{10}{100}\right)$ & 
\centering \vspace{-0.3cm} \textbf{C.} $200 \times \left(1+\dfrac{10}{100}\right)$ & 
\centering \vspace{-0.3cm} \textbf{D.} $200 \times 0,9$ \tabularnewline
\hline
\end{tabular}

%-------------------------
\vspace{0.7cm}

\textbf{Question 6 :} 
\vspace{5pt}

Un article coûte 400 € et subit une augmentation de 20\%.  
Son nouveau prix est : 
\vspace{0.3cm}

\renewcommand{\arraystretch}{3} 
\begin{tabular}{|m{3.4cm}|m{3.4cm}|m{3.6cm}|m{3.4cm}|}
\hline
\centering \vspace{-0.3cm} \textbf{A.} $420$ € & 
\centering \vspace{-0.3cm} \textbf{B.} $480$ € & 
\centering \vspace{-0.3cm} \textbf{C.} $500$ € & 
\centering \vspace{-0.3cm} \textbf{D.} $320$ € \tabularnewline
\hline
\end{tabular}

%-------------------------
\vspace{0.7cm}

\textbf{Question 7 :} 
\vspace{5pt}

Le prix d'un produit après augmentation est de 
120 € et l'augmentation était de 20\%.  
Le prix initial était : 
\vspace{0.3cm}

\renewcommand{\arraystretch}{4} 
\begin{tabular}{|m{3.4cm}|m{3.4cm}|m{3.4cm}|m{3.4cm}|}
\hline
\centering \vspace{-0.5cm} \textbf{A.} $120 \times 0,8$ & 
\centering \vspace{-0.9cm} \textbf{B.} $\dfrac{120}{1 - \dfrac{20}{100}}$ & 
\centering \vspace{-0.5cm} \textbf{C.} $\dfrac{120}{1,2}$ & 
\centering \vspace{-0.5cm} \textbf{D.} $120 : 0,2$ \tabularnewline
\hline
\end{tabular}

\vspace{0.7cm}

\textbf{Question 8 :} 
\vspace{5pt}

Après une diminution de 25\%, un article vaut 90 €.  
Sa valeur initiale était : 
\vspace{0.3cm}

\renewcommand{\arraystretch}{3} 
\begin{tabular}{|m{3.4cm}|m{3.4cm}|m{3.4cm}|m{3.4cm}|}
\hline
\centering \vspace{-0.3cm} \textbf{A.} 100 € & 
\centering \vspace{-0.3cm} \textbf{B.} 120 € & 
\centering \vspace{-0.3cm} \textbf{C.} 130 € & 
\centering \vspace{-0.3cm} \textbf{D.} 125 € \tabularnewline
\hline
\end{tabular}






\vspace{0.7cm}


\end{document}
